\chapter{从已有数据自己回测一天：完整演练}

\section{本章要解决的困惑}

前面每章都拆了一个对象。本章把它们重新装回一天的实验。目标不是记命令，而是让你知道每一步在验证哪条数据变换。

我们选一个单日实验：

\begin{lstlisting}
symbol = CCUSDT
date = 2026-05-18
strategy = old_tfi_four_cell_shadow_v1
admission = entry_spread_q70_v1
capacity = core_idle01 or fifo_clip
exit = fixed60_taker or fixed60_mid
\end{lstlisting}

如果你能理解这个单日实验，三天、多 profile、strict audit 都只是规模扩展。

\section{第一步：确认市场事实存在}

先确认 canonical 数据大致存在。你不需要先关心策略，只要知道有三条事实流：

\begin{lstlisting}
quote_frame_v1
trade_event_v1
l2_level_update_v1
\end{lstlisting}

它们分别回答：

\begin{itemize}[leftmargin=2em]
  \item 当前最优 bid/ask 是多少；
  \item 最近主动成交方向和数量是什么；
  \item 订单簿价位和数量怎样更新。
\end{itemize}

如果这三条事实流不完整，后面的策略结果都要标记证据不完整。尤其是 TFI 依赖 trade，不能把 quote-only 结果伪装成 TFI 完整证据。

\section{第二步：确认 decision frame 是 runtime-safe}

\code{decision_frame_v1} 是 fast 回测的热路径。它应该只包含 entry 时刻以前可见的信息。你检查它不是为了相信它，而是为了确认它没有混进未来答案。

合法字段包括：

\begin{lstlisting}
best_bid_price
best_ask_price
mid_price
spread_bps
trade_window_count
trade_flow_imbalance
frames_since_mid_change
past_event_25_bps
\end{lstlisting}

非法直觉包括：

\begin{lstlisting}
future_return_60s
mfe_after_entry
mae_after_entry
realized_pnl
scored_entry_quality
\end{lstlisting}

\warningline{如果某个策略输入字段只有 entry 之后才能知道，它就不能进 runtime。}

\section{第三步：拿一行 frame 做 entry 判断}

假设一行 frame：

\begin{lstlisting}
trade_flow_imbalance = -1.0
trade_window_count = 2
past_event_25_bps = -0.1
frames_since_mid_change = 38
spread_bps = 18.0
\end{lstlisting}

先判断方向：

\[
TFI=-1 \Rightarrow side=sell,\quad direction=-1.
\]

再判断 flat：

\[
|-0.1|\le0.5 \Rightarrow flat=true.
\]

再判断 stale：

\[
38\ge25 \Rightarrow stale25=true,\quad 38\ge31 \Rightarrow B=true.
\]

如果当天 prior-day q70 entry cross threshold 是 10 bps：

\[
entry\_cross=\frac{18.0}{2}=9.0\le10.
\]

admission 通过。

\section{第四步：算 cell 和仓位}

假设闭合记忆：

\begin{lstlisting}
closed_bps = [1.0, 0.5, -0.4, 0.8, -0.2]
\end{lstlisting}

\[
P_5=1.0+0.5+0.8=2.3,
\quad
N_5=0.4+0.2=0.6.
\]

\[
R5=\frac{2.3}{0.6}=3.83.
\]

于是：

\[
A=1,\quad B=1,\quad \text{cell}=\text{\ttfamily 11\_r5\_frames}.
\]

假设 membership set 是：

\begin{lstlisting}
tfi_follow_flat+tfi_short_flat+tfi_short_stale25+tfi_event_active
\end{lstlisting}

则：

\[
base\_weight=2.0.
\]

frames 没超过 59.8：

\[
weak\_overlay\_weight=2.0.
\]

默认 gamma：

\[
\gamma_{11}=0.75.
\]

目标仓位：

\[
target=2.0\times0.75=1.5.
\]

如果入场前 open exposure 是 2.0：

\[
actual=\min(1.5,3.0-2.0)=1.0.
\]

\section{第五步：算 fixed60 taker 收益}

short 入场时：

\begin{lstlisting}
entry_bid = 0.1120
entry_ask = 0.1124
\end{lstlisting}

60 秒后：

\begin{lstlisting}
exit_bid = 0.1114
exit_ask = 0.1118
\end{lstlisting}

short executable：

\[
Y^{exec,short}
=10000\log\frac{entry\_bid}{exit\_ask}
=10000\log\frac{0.1120}{0.1118}
\approx17.9\text{ bps}.
\]

如果 actual exposure 是 1.0，则这笔 weighted contribution 近似就是 17.9 bps exposure-units。若 actual exposure 只有 0.2，则贡献按 0.2 缩小。

\section{第六步：把手算对应到输出文件}

你应该能在 \code{entries.parquet} 找到：

\begin{lstlisting}
side = sell
cell = 11_r5_frames
membership_set = ...
entry_cross_bps = 9.0
admission_accepted = true
requested_exposure = 1.5
actual_exposure = 1.0
\end{lstlisting}

在 \code{exits.parquet} 找到：

\begin{lstlisting}
exit_profile = fixed60_taker
entry_bid = 0.1120
exit_ask = 0.1118
net_bps ~= 17.9
\end{lstlisting}

在 \code{daily.csv} 看到这笔贡献进入当天汇总。

\section{第七步：读 summary 不要只看一个数}

\code{summary.json} 最少看：

\begin{lstlisting}
policy
admission_profile
capacity_profile
exit_profile
actual_entries
admission_accepted_entries
admission_rejected_entries
exits
net_weighted_bps
daily rows
\end{lstlisting}

如果一个 summary 没写清楚 profile，就不能拿来和别的结果比较。

\section{第八步：什么时候这个单日实验算成功}

不是赚了就成功。学习意义上的成功是：

\begin{enumerate}[leftmargin=2em]
  \item 你能解释每个 entry 为什么触发；
  \item 你能解释每个 actual exposure 为什么被剪裁或没被剪裁；
  \item 你能解释 mid 和 executable 为什么不同；
  \item 你能指出这次实验有没有用到 future label；
  \item 你能说明它是 fast evidence，不是 strict exchange audit。
\end{enumerate}

\section{把一天扩展到多天}

多天不是换一种算法。只是重复：

\begin{lstlisting}
for day in date_range:
    load decision_frame cache
    start/reset day-local buckets
    preserve or seed runtime-safe closed state as configured
    scan frames
    write daily rows
\end{lstlisting}

真正需要小心的是：

\begin{itemize}[leftmargin=2em]
  \item q70 threshold 是否来自 prior day；
  \item R5 state 是否跨日按 runtime-safe 规则处理；
  \item quote/trade/L2 覆盖是否一致；
  \item 不同日期是否用了同一 profile；
  \item worst day 是否暴露结构问题。
\end{itemize}

\checkpoint{本章的目标是让你能闭着眼睛画出一条链：frame 到 signal，到 admission，到 capacity，到 entry，到 exit，到 daily。}
